%
% This work is licensed under a Creative Commons Attribution-ShareAlike 4.0 International License.
%

\documentclass[10pt, letterpaper, twoside, openany]{article}

\title{Introduction to Quantum Mechanics -- Worksheet 6}
\author{Joseph D. MacMillan}
\date{}


\usepackage{graphicx}
\usepackage{color} 
\usepackage{amsmath, amssymb}
\usepackage{libertinus}
\usepackage{microtype}
\usepackage{layout}
\usepackage[most]{tcolorbox}
\tcbuselibrary{skins,breakable}
\usepackage[hang, small, bf]{caption}
\captionsetup[table]{position=top}
\captionsetup[figure]{position=bottom}
\usepackage[hyphens]{url}
\usepackage[hidelinks]{hyperref}
\hypersetup{pdftitle={Introduction to Quantum Mechanics -- Worksheet 6}}
\hypersetup{pdfauthor={Joseph D. MacMillan}}
\urlstyle{rm}
\usepackage[shortlabels]{enumitem}
\usepackage[margin=1in]{geometry}
\usepackage{fancyhdr}




\newcommand{\uu}{\symbfup{u}}
\newcommand{\grad}{\symbfup{\nabla}}
\newcommand{\curl}{\symbfup{\nabla} \times}
\newcommand{\dfdx}[2]{\frac{\partial {#1}}{\partial {#2}}}
\newcommand{\ddfdx}[2]{\frac{\partial^2 {#1}}{\partial {#2}^2}}
\newcommand{\vort}{\symbfup{\omega}}
\renewcommand\vec{\symbfup}
\newcommand{\unit}[1]{\hat{\vec{#1}}}
\newcommand{\U}{\mathbb{U}}
\newcommand{\V}{\mathbb{V}}
\newcommand{\LL}{\mathbb{L}^2}
\newcommand{\LLLL}{\mathbb{L}^4}

\newcommand{\ket}[1]{| #1 \rangle}
\newcommand{\bra}[1]{\langle #1 |}
\newcommand{\braket}[2]{\langle #1 | #2 \rangle}
\newcommand{\avg}[1]{\langle #1 \rangle}


%\DeclareMathAlphabet{\mathbfsf}{\encodingdefault}{\sfdefault}{bx}{sl}
\newcommand{\tens}[1]{\mathbfsfit{#1}}


\newcounter{example}[section]
\def\theexample{\thechapter.\arabic{example}}
\newenvironment{example}[1][ ]{\refstepcounter{example}
\begin{tcolorbox}[breakable, sharp corners, boxrule = 0pt, frame empty, opacityframe=0, parbox=false]
\textbf{Example \theexample \ -- #1.}
}
{ 
\end{tcolorbox}
}

\newenvironment{theorem}[1][ ]{
\begin{tcolorbox}[breakable, sharp corners, boxrule = 0pt, frame empty, opacityframe=0, parbox=false]
\textbf{#1.}
}
{ 
\end{tcolorbox}
}

\newcounter{problem}[section]
\def\theproblem{\thechapter.\arabic{problem}}
\newenvironment{problem}[1][ ]{\refstepcounter{problem} \noindent \textbf{Problem \theproblem \ -- #1.}}{\vspace{0.1in}}

\renewcommand{\arraystretch}{2.5}

\begin{document}

\pagestyle{fancy}
\fancyhead[L]{\emph{Introduction to Quantum Mechanics}}
\fancyhead[R]{Worksheet 6} 

\Large \noindent Worksheet 6

\noindent \LARGE \textbf{The Infinite Square Well}

\normalsize

\vspace{0.2in}

\noindent Position space  \bullet \  Allowed energies   \bullet \ Energy eigenstates \bullet \ Potential well  

\vspace{0.1in}

\noindent \rule{\textwidth}{0.5pt}

\vspace{0.2in}

\begin{enumerate}

\item The infinite square well -- also called a particle in a box -- is defined by the potential energy
\[
V(x) = 
\begin{cases} 
\infty, &  x < 0 \\
0, & 0 < x < a \\
\infty, & x>a,
\end{cases}
\]
I want you to draw the axes for two plots, one for $\phi_n(x)$ and one for $|\phi_n(x)|^2$.  Make the horizontal axes a little up from the bottom of the page, and draw two vertical axes -- one to the far left and one more in the middle.  Label the left axis $\phi_n(x)$, the right axis $|\phi_n(x)|^2$, and the horizontal axes $x$.  On the $x$ axes, label the width $a$.  Finally, sketch on the potential energy as a dark, thick line on each plot.

\newpage

\item Think about this carefully: if the particle is \emph{classical}, what is the probability of finding it near some position $x$?  For example, what's the probability of finding it on the left-hand-side of the box, or the middle 25\% of the box?  Write out an equation for the probability density $\mathcal{P}(x)$; don't forget to properly normalize it.  Then, most importantly, sketch it on the $|\phi_n(x)|^2$ graph on the previous page, using a faint thin line.

\vspace{2in}

\item Our job today is to solve the energy eigenvalue equation, 
\[
-\frac{\hbar^2}{2m} \frac{d^2 \phi_n}{dx^2} + V(x)\phi_n(x) = E_n \phi_n(x)
\]
for the wave functions $\phi_n(x)$ and allowed energies $E_n$.  But for most of the region $-\infty < x < \infty$, that's easy:  what should the wave function $\phi(x)$ be \emph{outside} the well?  Write your answer in the form of a \emph{boundary condition}.

\vspace{2in}

\item Now we're ready to solve the differential equation for the region $0<x<a$.  Write it out for that region, and then ``make it pretty'' by defining a new constant to hold all your other constants and rearranging for the highest derivative.

\newpage

\item This is a second order ODE, so we expect two independent functions as a solution.  If you know the solutions already, great!  Write out the general solution below.  If not -- guess and check first!

\vspace{2in}

\item You should now have a solution with three unknown constants.  What are they?  Use your boundary condition at $x=0$ to find one of your constants.

\vspace{2in}

\item Use your boundary condition at $x=a$ to find a second constant.  This one is tricky, so go carefully and remember what you know about the sine function.  How many possibilities can this constant be?

\newpage

\item Write out the allowed energies of the system in terms of the width $a$, the particle's mass $m$, and Planck's constant $\hbar$, and then write out your wave function $\phi(x)$ using everything that you know so far.

\vspace{2in}

\item You still have one constant (mine was $A$) left -- how do you figure out that one?  Do it, and write your final wave function $\phi_n(x)$.

\vspace{3in}

\item Sketch the first three energy eigenstates on your plots.  But -- and sorry, this is a little confusing -- don't use your $x$-axis, draw a new one for each plot that (relatively) matches the energy of each state.  I'll explain better in person if you like.

\vspace{1in}

\item Finally, sketch (as best you can, things are getting crowded) the shape of $|\phi_n(x)|^2$ for $n=20$.  Compare your quantum probability densities with the classical one; any conclusions?

\end{enumerate}

\end{document}
